\chapter{Introducción y alcance}
El problema del ciclo hamiltoniano consiste en determinar si un grafo contiene un ciclo que visite todos sus vértices exactamente una vez antes de regresar al inicio. Su dificultad combinatoria motiva el estudio de métodos de búsqueda heurística. Esta investigación corresponde a la fase previa a la implementación: fundamenta la representación, la evaluación y las decisiones de diseño de SA, ACO, colonia de abejas, DE y PSO. No presenta resultados experimentales.

El documento integra las investigaciones de las ramas del equipo distintas de main y las complementa con el archivo local investigacion.tex. Se conservan los ejemplos y los pseudocódigos de las aportaciones; los parámetros y adaptaciones se identifican como propuestas, no como resultados comprobados sobre FHCP. Bee Colony se concreta en la variante Artificial Bee Colony (ABC), para mantener una descripción inequívoca de sus fases.

El texto se organiza en cinco bloques: grafos y SA; representación y ACO; heurísticas y ABC; métricas y DE; y FHCP y PSO. La síntesis final relaciona estos bloques con un protocolo de comparación. Las referencias heredadas conservan las notas sobre el alcance de su consulta; una referencia que el autor original declaró pendiente no se presenta como verificada íntegramente en esta integración.
